% SOURCE: AI JMP.pdf p. 89, Table G.1. Transcribed, not regenerated.
\begin{table}[H]
\centering
\caption{Equity-beta benchmark sensitivity}
\label{tab:beta-benchmark-sensitivity}
\small
\begin{tabular}{llrrrr}
\toprule
Sample & Equity benchmark & Spot effect & Wild $p$ & Total effect & Wild $p$ \\
\midrule
Maximal all-nine-beta sample & S\&P 500 & $-7.93$ & 0.005 & $-7.19$ & 0.009 \\
Maximal all-nine-beta sample & MSCI World & $-2.93$ & 0.319 & $-4.00$ & 0.162 \\
Post-1999 & S\&P 500 & $-11.10$ & 0.003 & $-10.73$ & 0.004 \\
Post-1999 & MSCI World & $-7.67$ & 0.030 & $-7.85$ & 0.019 \\
\bottomrule
\end{tabular}

\begin{minipage}{0.94\textwidth}
\footnotesize\emph{Notes:} Annualized percentage-point state coefficients for portfolios sorted on expanding equity betas. The full-sample MSCI result is weak. The original explanation is home-index contamination from Japan's historically large MSCI World weight; after 1999 the two benchmarks agree in sign and statistical detection. The table shows that benchmark construction is a material part of the beta robustness claim.
\end{minipage}
\end{table}
